% !TEX TS-program = xelatex
% !TEX encoding = UTF-8 Unicode
% AUTO-GENERATED by scripts/build-latex-cv.mjs from canonical site data (src/data/*).
% Do not edit manually; run `npm run build:latex` after updating src/data/.
\documentclass[10pt,a4paper]{article}
\usepackage[a4paper,left=19mm,right=19mm,top=16mm,bottom=18mm,footskip=9mm]{geometry}
\usepackage{iftex}
\ifTUTeX
  \usepackage{fontspec}
  \usepackage{xeCJK}
  \IfFileExists{/System/Library/Fonts/STHeiti Light.ttc}{%
    \setCJKmainfont{STHeiti Light.ttc}[Path=/System/Library/Fonts/,FontIndex=0,BoldFont=STHeiti Medium.ttc]
    \setCJKsansfont{STHeiti Light.ttc}[Path=/System/Library/Fonts/,FontIndex=0,BoldFont=STHeiti Medium.ttc]
  }{%
    \setCJKmainfont{FandolHei-Regular.otf}[BoldFont=FandolHei-Bold.otf]
    \setCJKsansfont{FandolHei-Regular.otf}[BoldFont=FandolHei-Bold.otf]
  }
  \XeTeXlinebreaklocale "zh"
  \XeTeXlinebreakskip=0pt plus 0.5pt
\else
  \usepackage[UTF8]{ctex}
\fi
\usepackage{xcolor}
\definecolor{cvaccent}{HTML}{0E7490}
\definecolor{cvmuted}{HTML}{475569}
\usepackage{enumitem}
\usepackage{needspace}
\usepackage{fancyhdr}
\usepackage{lastpage}
\usepackage[unicode,colorlinks=true,urlcolor=cvaccent,linkcolor=cvaccent,pdfborder={0 0 0}]{hyperref}
\hypersetup{
  pdftitle={简历 — 李栋林（Dennis）},
  pdfauthor={李栋林（Dennis）},
  pdfcreator={scripts/build-latex-cv.mjs}
}
\urlstyle{same}
\setlength{\parindent}{0pt}
\setlength{\parskip}{0pt}
\setlength{\emergencystretch}{2em}
\setlength{\headheight}{13pt}
\setlength{\headsep}{3mm}
\renewcommand{\baselinestretch}{1.15}
\raggedbottom
\widowpenalty=10000
\clubpenalty=10000
\displaywidowpenalty=10000
\setlist[itemize]{leftmargin=4.5mm,label=\textbullet,itemsep=2pt,topsep=2pt,parsep=0pt,partopsep=0pt}
\pagestyle{fancy}
\fancyhf{}
\renewcommand{\headrulewidth}{0pt}
\renewcommand{\footrulewidth}{0pt}
\fancyfoot[L]{\color{cvmuted}\fontsize{8.3}{10}\selectfont\href{https://donglinli.com/zh/cv/}{donglinli.com/zh/cv}}
\fancyfoot[C]{\color{cvmuted}\fontsize{8.3}{10}\selectfont 简历 \textbar{} 更新于 2026 年 11 月 22 日}
\fancyfoot[R]{\color{cvmuted}\fontsize{8.3}{10}\selectfont\thepage\,/\,\pageref*{LastPage}}
\newcommand{\CVHeader}[3]{%
  {\fontsize{20}{24}\selectfont\bfseries #1\par}%
  \vspace{2pt}{\fontsize{10.5}{13}\selectfont\color{cvmuted}#2\par}%
  \vspace{4pt}{\fontsize{9.2}{11.5}\selectfont #3\par}%
  \vspace{6pt}{\color{cvaccent}\hrule height 0.8pt}\vspace{6pt}}
\newcommand{\CVSection}[2]{%
  \Needspace{6\baselineskip}%
  \vspace{8pt}\phantomsection\pdfbookmark[1]{#2}{sec-#1}%
  {\fontsize{12}{15}\selectfont\bfseries\color{cvaccent} #2\par}%
  \vspace{2pt}{\color{cvaccent}\hrule height 0.35pt}\vspace{5pt}}
\newcommand{\CVEntry}[3]{%
  \Needspace{4\baselineskip}%
  \vspace{2pt}\noindent\begin{tabular*}{\textwidth}{@{\extracolsep{\fill}}l r@{}}%
    \parbox[t]{0.76\textwidth}{\raggedright\textbf{#1}} & {\color{cvmuted}\fontsize{9.4}{11.5}\selectfont #2}%
  \end{tabular*}\par
  \if\relax\detokenize{#3}\relax\else
    \vspace{1pt}{\color{cvmuted}\fontsize{9.6}{12}\selectfont #3\par}%
  \fi}
\newcommand{\CVText}[1]{\vspace{2pt}{\fontsize{9.6}{12.5}\selectfont #1\par}}
\newcommand{\CVMeta}[1]{\vspace{1.5pt}{\color{cvmuted}\fontsize{9.2}{11.5}\selectfont #1\par}}
\newcommand{\CVEntryEnd}{\vspace{4pt}}

\begin{document}
\CVHeader{李栋林（Dennis）}{曼彻斯特大学电气与电子工程博士研究生}{\href{mailto:hello@donglinli.com}{hello@donglinli.com} \textbar{} \href{https://donglinli.com/zh/}{donglinli.com/zh} \textbar{} \href{https://github.com/Daihaolin201}{GitHub} \textbar{} \href{https://www.linkedin.com/in/donglin-li-55424326b}{LinkedIn} \textbar{} \href{https://orcid.org/0009-0007-2130-4862}{ORCID} \textbar{} \href{https://scholar.google.com/citations?user=mBr54wEAAAAJ}{谷歌学术}}

\CVSection{profile}{个人简介}
\CVText{\textbf{让多机器人在不确定环境中协同决策。}}
\CVText{曼彻斯特大学电气与电子工程博士研究生，研究多机器人协同寻源的贝叶斯对偶控制，同时开发嵌入式与自主机器人系统和科研软件；在校讲授控制、机器人与数字电子方向的实验课和习题课。}

\CVSection{education}{教育}
\CVEntry{\href{https://www.manchester.ac.uk/study/postgraduate-research/programmes/list/04620/phd-electrical-and-electronic-engineering/}{电气与电子工程博士（在读）}}{2024.09—至今}{曼彻斯特大学}
\CVEntryEnd
\CVEntry{电子工程荣誉学士（一等荣誉学位）}{2021.09—2024.06}{曼彻斯特大学}
\CVEntryEnd

\CVSection{research}{研究}
\CVEntry{多机器人协同寻源的贝叶斯对偶控制}{2024.09—至今}{曼彻斯特大学 · 博士研究生}
\begin{itemize}
  \item 面向未知含噪场中的多机器人协同寻源，研究贝叶斯对偶控制：每台机器人的运动同时兼顾降低源位置估计的不确定性与向估计位置收敛。
  \item 提出基于一致性均值融合与安全投影的对偶控制方法（CCC 2026），以及面向通信效率的迹加权事件触发均值融合（IAI 2026，最佳论文奖入围）。
\end{itemize}
\CVEntryEnd

\CVSection{publications}{论文}
\begin{enumerate}[leftmargin=5.5mm,itemsep=4pt,topsep=2pt,parsep=0pt]
  \item \textbf{Donglin Li}, Zhongguo Li, Hanbo Ma, Chunyi Cao, Zhengtao Ding。\href{https://cms.amss.ac.cn/conferences/ccc26/program/TuSessions.html\#TuA11-4}{Dual Control for Multi-Robot Source Seeking with Consensus Mean Fusion and Safety Projection}。第45届中国控制会议（CCC 2026），中国长沙，2026。\\
  {\color{cvmuted}\fontsize{9.4}{11.5}\selectfont 第一作者 \textbar{} 已接收。}
  \item \textbf{Donglin Li}, Zhongguo Li, Junan Wang, Hanbo Ma, Zhengtao Ding。Trace-Weighted Event-Triggered Mean Fusion for Efficient Multi-Robot Source Seeking。第8届工业人工智能国际会议（IAI 2026），中国沈阳，2026。\\
  {\color{cvmuted}\fontsize{9.4}{11.5}\selectfont 第一作者 \textbar{} 已接收；最佳论文奖入围。}
  \item Hanbo Ma, Zhongguo Li, Junan Wang, Chunyi Cao, \textbf{Donglin Li}, Jun Yang, Zhengtao Ding。\href{https://doi.org/10.1016/j.neucom.2026.134973}{Load Balancing for Large Language Model Inference Serving: A Control and Learning Oriented Survey}。Neurocomputing 706（2026），134973。\\
  {\color{cvmuted}\fontsize{9.4}{11.5}\selectfont 共同作者 \textbar{} 2026 年 9 月 3 日在线发表。 \textbar{} \href{https://doi.org/10.1016/j.neucom.2026.134973}{DOI: 10.1016/j.neucom.2026.134973}}
  \item Hanbo Ma, Zhongguo Li, Chunyi Cao, \textbf{Donglin Li}, Zhengtao Ding。\href{https://ifac-sms.org/conferences/conferences/IFAC26/program/IFAC26_ContentListWeb_4.html\#thc04_03}{Multi-Objective Pipeline Scheduling for LLM Inference in Heterogeneous DePIN}。第23届 IFAC 世界大会（IFAC WC 2026），韩国釜山，2026。\\
  {\color{cvmuted}\fontsize{9.4}{11.5}\selectfont 共同作者 \textbar{} 已接收。}
\end{enumerate}

\CVSection{employment}{经历}
\CVEntry{研究生助教}{2024.09—至今}{曼彻斯特大学}
\begin{itemize}
  \item 承担控制、机器人与数字电子方向六门本科及研究生课程的实验与习题课教学，指导控制器实现、机器人运动学和电路设计。
  \item 负责作业与考试评阅及书面反馈，为毕业设计展示提供技术评议；参与 2026 年夏季阅卷与课程开发，并获聘担任 2026/27 学年 EEEN11101 课程助教。
\end{itemize}
\CVEntryEnd
\CVEntry{客座讲师与实验设计者（过程控制）}{2025.06—2025.12}{曼彻斯特大学}
\begin{itemize}
  \item 设计并搭建硕士过程控制实验平台：PLC 执行 PID 控制，经 OPC 与 MATLAB/Simulink 互联，实现实时监控与分析。
  \item 承担相关控制理论的客座讲授，并主持实验教学。
\end{itemize}
\CVEntryEnd

\CVSection{projects}{项目}
\CVEntry{\href{https://hackabot-2026.com/live/project-4/}{多机器人协同推车}}{2026.03}{Hack-a-Bot 2026，曼彻斯特大学 RoboSoc \textbar{} 三等奖 \textbar{} 团队负责人}
\CVText{最多三台 MONA 差速机器人须绕开障碍、把推车推入目标区；感知只有一台共享的头顶相机，通信只有 Wi-Fi。感知、规划与编队控制集中运行在笔记本上，构成一个闭环。}
\CVText{\textbf{负责：}带队，并负责整套软件：ArUco 定位、留足安全距离的 A* 规划、Pure Pursuit 与 PID 跟踪，以及和机器人的通信}
\CVMeta{\textbf{技术栈：}OpenCV · ArUco · A* · Pure Pursuit · PID · TCP}
\CVEntryEnd
\CVEntry{\href{https://uom-robosoc.github.io/tasks/Task\%202.html}{MONA 自主配送机器人}}{2025.03}{Hack-a-Bot 2025，曼彻斯特大学 RoboSoc 与 Sony \textbar{} 一等奖}
\CVText{24 小时挑战：MONA 机器人须在一次限时运行中自主穿越胶带迷宫，完成起点 \ensuremath{\rightarrow} 取货点 \ensuremath{\rightarrow} 卸货点的配送，定位只依赖头顶摄像头的视频流。}
\CVText{\textbf{负责：}头顶视觉、A* 规划和 PID 控制，以及无线链路和电机校准的 Arduino 代码}
\CVMeta{\textbf{技术栈：}OpenCV · ArUco · A* · PID · NRF24L01 · Arduino}
\CVEntryEnd
\CVEntry{\href{https://github.com/Daihaolin201/operator-one}{OperatorOne}}{2026.03}{UK AI Agent Hackathon EP4 \ensuremath{\times} OpenClaw，伦敦帝国理工学院}
\CVText{一个把商业想法推进到产品、市场、销售与运营的多智能体系统：CEO 编排器以分阶段门控的循环调度四个 OpenClaw 专职智能体，向收入目标迭代。}
\CVText{\textbf{负责：}CEO 编排脚本，仪表板的运行接口和时间线，GLM-5 模型接入和预检，真实外发动作的安全开关，以及参赛提交材料}
\CVMeta{\textbf{技术栈：}OpenClaw · Python · GLM-5 · JSON · GitHub Actions}
\CVEntryEnd
\CVEntry{循迹与交通标志识别}{2023.04}{NVIDIA PuzzleThon，Manchester Robotics 与曼彻斯特大学 \textbar{} 二等奖}
\CVText{自主 PuzzleBot（Jetson Nano、ROS）须在未知赛道上无人干预地行驶：循线、遵守红绿灯、按方向标志动作。}
\CVText{\textbf{负责：}颜色滤波与掩膜、循迹和标志识别}
\CVMeta{\textbf{技术栈：}ROS · Python · OpenCV · Jetson Nano}
\CVEntryEnd
\CVEntry{\href{https://github.com/Daihaolin201/Smart-TrafficCones}{智能交通锥}}{2023.09—2024.05}{本科三年级个人项目，曼彻斯特大学}
\CVText{可自主布设的机器人交通锥：四台锥桶车（一台领航、三台跟随）按设定间距列队并复现领航车的路线，由手持控制器经 802.15.4 无线网络指挥。}
\CVText{\textbf{负责：}ESP32-S3 控制板、ESP32 和 STM32 固件、无线协议、控制回路和实验，全部独立完成}
\CVMeta{\textbf{技术栈：}ESP32-S3 · STM32 · XBee 802.15.4 · PCB · PID}
\CVEntryEnd
\CVEntry{ESP 自主循迹小车}{2022.09—2023.05}{嵌入式系统项目，曼彻斯特大学 \textbar{} 三等奖}
\CVText{基于 STM32F401（Mbed C++）的自主循迹小车，赛道包含坡道、180\ensuremath{^\circ} 掉头与受控停车；五人团队从自制传感器板做起。}
\CVText{\textbf{负责：}电路设计和传感器板（Altium 迭代三版）、软件，以及期末报告的系统部分}
\CVMeta{\textbf{技术栈：}STM32 · Mbed C++ · Altium PCB · PID · BLE}
\CVEntryEnd
\CVEntry{CSF Pilot 与科研工具}{2026.07—至今}{桌面软件：集群作业、文献、论文图}
\CVText{CSF Pilot 是曼彻斯特大学 CSF3 高性能计算集群的桌面客户端（PySide6）。它维持一条经多因素认证的 SSH 会话，在其上复用文件传输、Slurm 作业管理、编辑器、终端和用量分析，并通过 MCP 服务把这条会话开放给 AI 智能体。}
\CVMeta{\textbf{技术栈：}PySide6 · paramiko · Slurm · MCP · pytest}
\CVEntryEnd
\CVEntry{教学工作台}{2026.09—至今}{家教与学生服务用的本地浏览器应用}
\CVText{本地优先的单文件 Web 应用，覆盖家教业务全流程：学生、排课、备课、作业、课后反馈、课时包与工资对账。权威数据是用户所选文件夹中的一个 JSON 文件，经 File System Access API 读写。}
\CVMeta{\textbf{技术栈：}HTML · JavaScript · File System Access API · ICS · node:test}
\CVEntryEnd

\CVSection{awards}{荣誉}
\begin{itemize}[leftmargin=4.5mm,label=\textbullet,itemsep=2.5pt,topsep=2pt,parsep=0pt]
  \item \textbf{2026} 第8届工业人工智能国际会议（IAI 2026）最佳论文奖入围
  \item \textbf{2026} Hack-a-Bot 机器人挑战赛三等奖 {\color{cvmuted}--- 曼彻斯特大学 RoboSoc}
  \item \textbf{2026} \href{https://www.se.manchester.ac.uk/doctoral-academy/support/finance-and-funding/doctoral-training-and-development-support-fund/}{博士培训与发展支持基金（DTDSF）} {\color{cvmuted}--- 曼彻斯特大学}
  \item \textbf{2025} 电气与电子工程研究生研究会议海报一等奖 {\color{cvmuted}--- 曼彻斯特大学}
  \item \textbf{2025} Hack-a-Bot 机器人挑战赛一等奖 {\color{cvmuted}--- 曼彻斯特大学 RoboSoc 与 Sony}
  \item \textbf{2023} NVIDIA PuzzleThon 二等奖
  \item \textbf{2023} ESP 自主循迹小车开发三等奖
\end{itemize}

\CVSection{service}{服务}
\CVEntry{学术论文审稿人}{2026.05—2026.10}{TIMC · IEEE TAES · JACA 2026}
\begin{itemize}
  \item 为 IEEE Transactions on Aerospace and Electronic Systems 和 TIMC 各审稿 1 篇，为 JACA 2026 审稿 3 篇。
\end{itemize}
\CVEntryEnd
\CVEntry{EEE 博士生代表}{2024.10—至今}{曼彻斯特大学电气与电子工程系}
\begin{itemize}
  \item 2025/26 学年牵头本系博士生代表工作，汇总博士生与助教的意见，在院系和 Doctoral Academy 层面代为反映并跟进。
\end{itemize}
\CVEntryEnd
\CVEntry{学生大使兼博士生大使（PGR Ambassador）}{2024.10—至今}{曼彻斯特大学}
\begin{itemize}
  \item 在开放日和外展活动中代表学校介绍工程专业与校园生活，并在实验室主持机器人项目演示。
\end{itemize}
\CVEntryEnd

\CVSection{skills}{技能}
\begin{itemize}[leftmargin=4.5mm,label=\textbullet,itemsep=2.5pt,topsep=2pt,parsep=0pt]
  \item \textbf{研究方法：} 贝叶斯估计与对偶控制；实验设计；蒙特卡洛仿真与消融实验；统计分析与优化；文献综述、学术写作和同行评审。
  \item \textbf{编程：} C/C++；Python（NumPy、SciPy、pandas、Matplotlib）；MATLAB；R；JavaScript/TypeScript；SQL；PIC18 汇编；PLC 梯形图/SCL；Bash；HTML/CSS；Linux。
  \item \textbf{控制与机器人：} PID、分布式与一致性控制、编队与主动源搜索；A* 和 Pure Pursuit；多机器人协同；ROS、OpenCV/ArUco 视觉定位。
  \item \textbf{建模与仿真：} MATLAB/Simulink；LabVIEW/NI myDAQ；PLECS；Multisim；Tanner T-Spice/SPICE；ANSYS。
  \item \textbf{硬件与嵌入式：} STM32/Mbed、PIC18F8722、ESP32/ESP8266、Arduino、Xilinx FPGA/VHDL；PCB 原型；裸机与 RTOS 固件。
  \item \textbf{教学与评价：} 实验与习题课教学；控制和机器人实验演示；作业与考试评阅；学生项目技术反馈。
\end{itemize}

\end{document}
